\documentclass[11pt]{article}
\input{shared/project_style.tex}
\lhead{MHD\_BoundaryCompression}
\rhead{Stage 6 Notes}
\begin{document}
\DocTitle{Implementation Notes -- Stage 6}{Curated baseline study suite}{Purpose: explain how the repository moved from generic sweeps to a small named comparison suite whose runs each test a distinct physical question about amplitude, pulse duration, or symmetry breaking.}
\begin{overviewbox}
\textbf{Document purpose.} Purpose: explain how the repository moved from generic sweeps to a small named comparison suite whose runs each test a distinct physical question about amplitude, pulse duration, or symmetry breaking.
\end{overviewbox}
\section{Symbol and acronym table}
\begin{symtable}
suite & curated run set & Small named collection of comparison cases with clear scientific roles. \\
weak / moderate / strong & drive-amplitude labels & Comparative levels of inward boundary forcing. \\
shorter / longer pulse & pulse-duration labels & Comparative changes in temporal drive width. \\
asymmetric timing perturbation & symmetry-breaking case & Case used to probe sensitivity to imperfect synchronization. \\
$J$ & objective function & Scalar used to rank the curated cases under the upgraded diagnostic framework. \\
\end{symtable}

\section{Purpose of Stage 6}
A generic parameter sweep is useful for broad scanning, but it is not always the best vehicle for scientific explanation. Stage 6 therefore introduces a curated study suite: a small set of named runs chosen because each one tests a specific physical idea.

\section{The curated run set}
The current suite contains six baseline runs:
\begin{enumerate}[leftmargin=1.8em]
    \item weak symmetric pulse,
    \item moderate symmetric pulse,
    \item strong symmetric pulse,
    \item shorter symmetric pulse,
    \item longer symmetric pulse,
    \item asymmetric timing perturbation.
\end{enumerate}
This set is intentionally compact. The aim is to produce a first interpretable figure and table set rather than an exhaustive optimization campaign.

\section{Why this is better than a blind sweep}
Each curated case has a clear explanatory role. The amplitude trio probes how strongly the plasma can be driven before uniformity begins to degrade. The pulse-width variations probe temporal shaping. The timing-perturbation case probes sensitivity to loss of synchronization and therefore to symmetry breaking.

\section{Recommended interpretation order}
The outputs of the curated suite should be read in the following order:
\begin{enumerate}[leftmargin=1.8em]
    \item compare peak objective values,
    \item compare compression ratios at the best times,
    \item compare nonuniformity metrics,
    \item inspect symmetry and stagnation penalties,
    \item check divergence behavior before trusting the ranking.
\end{enumerate}
This order avoids the mistake of ranking runs by one dramatic scalar while ignoring whether the compressed state remained scientifically credible.

\section{Why this stage matters}
Stage 6 is the point where the repository starts to behave like a real study rather than a solver-development exercise. The code can now generate a compact, named, discussion-ready comparison set that can motivate the next controlled family of studies, such as spatial tapering or pulse trains.

\end{document}
